% --- Archon physics formalization source begin ---
% archon:physics
% archon:covers PhyXMiniProblems/problem_phyx_mini_0278.lean
% archon:source-report reports/phyx_mini/problem_phyx_mini_0278.source.json

\chapter{Physics problem phyx\_mini\_0278}
\label{ch:PhyXMiniProblems_problem_phyx_mini_0278}

\paragraph{Problem source.}
\#\# Physical scenario

A torsion pendulum consists of a metal disk with a wire running through its center and soldered in place. The wire is mounted vertically on clamps and pulled taut. Figure gives the magnitude \$\textbackslash{}tau\$ of the torque needed to rotate the disk about its center (and thus twist the wire) versus the rotation angle \$\textbackslash{}theta\$. The vertical axis scale is set by \$\textbackslash{}tau\_s = 4.0 \textbackslash{}times 10\textasciicircum{}\{-3\}\$ N\$\textbackslash{}cdot\$m. The disk is rotated to \$\textbackslash{}theta = 0.200\$ rad and then released. Figure also shows the resulting oscillation in terms of angular position \$\textbackslash{}theta\$ versus time \$t\$. The horizontal axis scale is set by \$t\_s = 0.40\$ s.

\#\# Question

What is the maximum angular speed \$d\textbackslash{}theta/dt\$ of the disk?

\#\# Answer choices

- A: A: 2.98 rad/s

- B: B: 3.02 rad/s

- C: C: 3.26 rad/s

- D: D: 3.14 rad/s

\#\# Figure caption (auxiliary; use the image as primary evidence)

The image contains two graphs.

\#\#\# Top Graph:
- **Axes**: 
  - X-axis labeled \textbackslash{}( \textbackslash{}theta \textbackslash{},(\textbackslash{}text\{rad\}) \textbackslash{}) ranging from 0 to 0.20.
  - Y-axis labeled \textbackslash{}( \textbackslash{}tau \textbackslash{}, (10\textasciicircum{}\{-3\} \textbackslash{}, \textbackslash{}text\{N\} \textbackslash{}cdot \textbackslash{}text\{m\}) \textbackslash{}) with a marker at \textbackslash{}( \textbackslash{}tau\_s \textbackslash{}).
  
- **Content**: 
  - A linear plot representing a straight line with a positive slope, indicating a direct relationship between \textbackslash{}( \textbackslash{}theta \textbackslash{}) and \textbackslash{}( \textbackslash{}tau \textbackslash{}).

\#\#\# Bottom Graph:
- **Axes**: 
  - X-axis labeled \textbackslash{}( t \textbackslash{}, (\textbackslash{}text\{s\}) \textbackslash{}) ranging from 0 to \textbackslash{}( t\_s \textbackslash{}).
  - Y-axis labeled \textbackslash{}( \textbackslash{}theta \textbackslash{}, (\textbackslash{}text\{rad\}) \textbackslash{}) ranging from -0.2 to 0.2.
  
- **Content**:
  - A sinusoidal wave indicating oscillatory behavior of \textbackslash{}( \textbackslash{}theta \textbackslash{}) over time \textbackslash{}( t \textbackslash{}).

\#\#\# Other Features:
- Grid lines present on both graphs.
- \textbackslash{}( t\_s \textbackslash{}) is marked on the time axis of the second graph.

\#\# Dataset metadata

Resonance and Harmonics; Physical Model Grounding Reasoning, Predictive Reasoning

\paragraph{Recorded answer/context.}
D: D: 3.14 rad/s

\paragraph{Figure/image path.}
/root/proposal\_for\_physic/science-mango/phyx\_mini\_run/phyx\_data/test\_image/278.png

\paragraph{Formalization target.}
create a compiling Lean file with sorry bodies at `PhyXMiniProblems/problem\_phyx\_mini\_0278.lean`. The Lean declarations must preserve the physical quantities, dimensions or dimensional roles, figure labels, governing-law hypotheses, and final relation expressed by this problem.
Use Mathlib/Physlib names found through LeanExplore where available. If a physics API is missing, introduce faithful local abstractions rather than scalar placeholder aliases.

\begin{theorem}[Physics formalization target]
\label{thm:physics:phyx_mini_0278:target}
\lean{PhyXMiniProblems.ProblemPhyXMini0278.problem_phyx_mini_0278}
\uses{def:physics:phyx-mini-0278:phyxminiproblems-problemphyxmini0278-angularvelocityinradianspersecond, def:physics:phyx-mini-0278:phyxminiproblems-problemphyxmini0278-torsionpendulumsetup, def:physics:phyx-mini-0278:phyxminiproblems-problemphyxmini0278-matchesproblemdata, def:physics:phyx-mini-0278:phyxminiproblems-problemphyxmini0278-matchessuppliedfigure, def:physics:phyx-mini-0278:phyxminiproblems-problemphyxmini0278-hasphysicaltorsionpendulumparameters, def:physics:phyx-mini-0278:phyxminiproblems-problemphyxmini0278-satisfiesidealtorsionpendulumlaws, def:physics:phyx-mini-0278:phyxminiproblems-problemphyxmini0278-recordeddatasetanswer, def:physics:phyx-mini-0278:phyxminiproblems-problemphyxmini0278-matchesanswerchoice}
The assigned autoformalize agent should translate this physics problem into Lean declarations in the covered file, with theorem and lemma proof bodies written as `by sorry`.
\end{theorem}
\begin{proof}
This is an autoformalization task, not a proof task. Produce statements that can later be proved by the physics prover without weakening the physical model.
\end{proof}

% --- Archon declaration topology begin ---
\section{Declaration topology}
\begin{definition}[Time Quantity]
  \label{def:physics:phyx-mini-0278:phyxminiproblems-problemphyxmini0278-timequantity}
  \lean{PhyXMiniProblems.ProblemPhyXMini0278.TimeQuantity}
  A nonnegative physical duration
\end{definition}
\begin{proof}
  This is the declared model definition of Time Quantity.
\end{proof}
\begin{definition}[Moment Of Inertia Quantity]
  \label{def:physics:phyx-mini-0278:phyxminiproblems-problemphyxmini0278-momentofinertiaquantity}
  \lean{PhyXMiniProblems.ProblemPhyXMini0278.MomentOfInertiaQuantity}
  A disk's moment of inertia, with dimension mass times length squared
\end{definition}
\begin{proof}
  This is the declared model definition of Moment Of Inertia Quantity.
\end{proof}
\begin{definition}[Torque Magnitude Quantity]
  \label{def:physics:phyx-mini-0278:phyxminiproblems-problemphyxmini0278-torquemagnitudequantity}
  \lean{PhyXMiniProblems.ProblemPhyXMini0278.TorqueMagnitudeQuantity}
  A nonnegative torque magnitude, measured coherently in newton-metres
\end{definition}
\begin{proof}
  This is the declared model definition of Torque Magnitude Quantity.
\end{proof}
\begin{definition}[Torsion Constant Quantity]
  \label{def:physics:phyx-mini-0278:phyxminiproblems-problemphyxmini0278-torsionconstantquantity}
  \lean{PhyXMiniProblems.ProblemPhyXMini0278.TorsionConstantQuantity}
  The torsion constant of the wire. A radian is dimensionless, so the physical dimension of N m / rad is mass times length squared divided by time squared. The separate name records its restoring-coefficient role
\end{definition}
\begin{proof}
  This is the declared model definition of Torsion Constant Quantity.
\end{proof}
\begin{definition}[Signed Angular Velocity Quantity]
  \label{def:physics:phyx-mini-0278:phyxminiproblems-problemphyxmini0278-signedangularvelocityquantity}
  \lean{PhyXMiniProblems.ProblemPhyXMini0278.SignedAngularVelocityQuantity}
  Signed angular velocity; radians are dimensionless, so its dimension is inverse time
\end{definition}
\begin{proof}
  This is the declared model definition of Signed Angular Velocity Quantity.
\end{proof}
\begin{definition}[time In Seconds]
  \label{def:physics:phyx-mini-0278:phyxminiproblems-problemphyxmini0278-timeinseconds}
  \lean{PhyXMiniProblems.ProblemPhyXMini0278.timeInSeconds}
  \uses{def:physics:phyx-mini-0278:phyxminiproblems-problemphyxmini0278-timequantity}
  Read a duration in seconds
\end{definition}
\begin{proof}
  This is the declared model definition of time In Seconds.
\end{proof}
\begin{definition}[moment Of Inertia In Kilogram Meter Squared]
  \label{def:physics:phyx-mini-0278:phyxminiproblems-problemphyxmini0278-momentofinertiainkilogrammetersquared}
  \lean{PhyXMiniProblems.ProblemPhyXMini0278.momentOfInertiaInKilogramMeterSquared}
  \uses{def:physics:phyx-mini-0278:phyxminiproblems-problemphyxmini0278-momentofinertiaquantity}
  Read a moment of inertia in SI kg m\textasciicircum{}2
\end{definition}
\begin{proof}
  This is the declared model definition of moment Of Inertia In Kilogram Meter Squared.
\end{proof}
\begin{definition}[torque Magnitude In Newton Meters]
  \label{def:physics:phyx-mini-0278:phyxminiproblems-problemphyxmini0278-torquemagnitudeinnewtonmeters}
  \lean{PhyXMiniProblems.ProblemPhyXMini0278.torqueMagnitudeInNewtonMeters}
  \uses{def:physics:phyx-mini-0278:phyxminiproblems-problemphyxmini0278-torquemagnitudequantity}
  Read a torque magnitude in SI newton-metres
\end{definition}
\begin{proof}
  This is the declared model definition of torque Magnitude In Newton Meters.
\end{proof}
\begin{definition}[torsion Constant In Newton Meters Per Radian]
  \label{def:physics:phyx-mini-0278:phyxminiproblems-problemphyxmini0278-torsionconstantinnewtonmetersperradian}
  \lean{PhyXMiniProblems.ProblemPhyXMini0278.torsionConstantInNewtonMetersPerRadian}
  \uses{def:physics:phyx-mini-0278:phyxminiproblems-problemphyxmini0278-torsionconstantquantity}
  Read a torsion constant in SI N m / rad
\end{definition}
\begin{proof}
  This is the declared model definition of torsion Constant In Newton Meters Per Radian.
\end{proof}
\begin{definition}[angular Velocity In Radians Per Second]
  \label{def:physics:phyx-mini-0278:phyxminiproblems-problemphyxmini0278-angularvelocityinradianspersecond}
  \lean{PhyXMiniProblems.ProblemPhyXMini0278.angularVelocityInRadiansPerSecond}
  \uses{def:physics:phyx-mini-0278:phyxminiproblems-problemphyxmini0278-signedangularvelocityquantity}
  Read a signed angular velocity in radians per second
\end{definition}
\begin{proof}
  This is the declared model definition of angular Velocity In Radians Per Second.
\end{proof}
\begin{definition}[Apparatus Component]
  \label{def:physics:phyx-mini-0278:phyxminiproblems-problemphyxmini0278-apparatuscomponent}
  \lean{PhyXMiniProblems.ProblemPhyXMini0278.ApparatusComponent}
  Mechanical components named in the source description
\end{definition}
\begin{proof}
  This is the declared model definition of Apparatus Component.
\end{proof}
\begin{definition}[Graph Panel]
  \label{def:physics:phyx-mini-0278:phyxminiproblems-problemphyxmini0278-graphpanel}
  \lean{PhyXMiniProblems.ProblemPhyXMini0278.GraphPanel}
  The two graph panels in the supplied bitmap
\end{definition}
\begin{proof}
  This is the declared model definition of Graph Panel.
\end{proof}
\begin{definition}[Axis Direction]
  \label{def:physics:phyx-mini-0278:phyxminiproblems-problemphyxmini0278-axisdirection}
  \lean{PhyXMiniProblems.ProblemPhyXMini0278.AxisDirection}
  Horizontal and vertical graph-axis roles
\end{definition}
\begin{proof}
  This is the declared model definition of Axis Direction.
\end{proof}
\begin{definition}[Figure Axis Label]
  \label{def:physics:phyx-mini-0278:phyxminiproblems-problemphyxmini0278-figureaxislabel}
  \lean{PhyXMiniProblems.ProblemPhyXMini0278.FigureAxisLabel}
  Physical labels, including units, printed on the graph axes
\end{definition}
\begin{proof}
  This is the declared model definition of Figure Axis Label.
\end{proof}
\begin{definition}[Curve Kind]
  \label{def:physics:phyx-mini-0278:phyxminiproblems-problemphyxmini0278-curvekind}
  \lean{PhyXMiniProblems.ProblemPhyXMini0278.CurveKind}
  Qualitative shapes of the two plotted traces
\end{definition}
\begin{proof}
  This is the declared model definition of Curve Kind.
\end{proof}
\begin{definition}[Torsion Pendulum Apparatus]
  \label{def:physics:phyx-mini-0278:phyxminiproblems-problemphyxmini0278-torsionpendulumapparatus}
  \lean{PhyXMiniProblems.ProblemPhyXMini0278.TorsionPendulumApparatus}
  \uses{def:physics:phyx-mini-0278:phyxminiproblems-problemphyxmini0278-apparatuscomponent}
  Verbal geometry of the clamped disk-and-wire torsion pendulum
\end{definition}
\begin{proof}
  This is the declared model definition of Torsion Pendulum Apparatus.
\end{proof}
\begin{definition}[Supplied Figure]
  \label{def:physics:phyx-mini-0278:phyxminiproblems-problemphyxmini0278-suppliedfigure}
  \lean{PhyXMiniProblems.ProblemPhyXMini0278.SuppliedFigure}
  \uses{def:physics:phyx-mini-0278:phyxminiproblems-problemphyxmini0278-timequantity, def:physics:phyx-mini-0278:phyxminiproblems-problemphyxmini0278-torquemagnitudequantity, def:physics:phyx-mini-0278:phyxminiproblems-problemphyxmini0278-graphpanel, def:physics:phyx-mini-0278:phyxminiproblems-problemphyxmini0278-axisdirection, def:physics:phyx-mini-0278:phyxminiproblems-problemphyxmini0278-figureaxislabel, def:physics:phyx-mini-0278:phyxminiproblems-problemphyxmini0278-curvekind}
  Data channels and labels belonging to the supplied two-panel figure. The functions store the plotted torque magnitudes and angular-position trace; their quantitative interpretation is given by readout assumptions below
\end{definition}
\begin{proof}
  This is the declared model definition of Supplied Figure.
\end{proof}
\begin{definition}[Torsion Pendulum Setup]
  \label{def:physics:phyx-mini-0278:phyxminiproblems-problemphyxmini0278-torsionpendulumsetup}
  \lean{PhyXMiniProblems.ProblemPhyXMini0278.TorsionPendulumSetup}
  \uses{def:physics:phyx-mini-0278:phyxminiproblems-problemphyxmini0278-momentofinertiaquantity, def:physics:phyx-mini-0278:phyxminiproblems-problemphyxmini0278-torsionconstantquantity, def:physics:phyx-mini-0278:phyxminiproblems-problemphyxmini0278-signedangularvelocityquantity, def:physics:phyx-mini-0278:phyxminiproblems-problemphyxmini0278-torsionpendulumapparatus, def:physics:phyx-mini-0278:phyxminiproblems-problemphyxmini0278-suppliedfigure}
  The torsion-pendulum experiment and its unknown time-dependent motion. The angular-position and angular-velocity functions are independent fields; neither is defined from the requested maximum speed. scalarOscillator is Physlib's one-coordinate harmonic oscillator, while amplitudePhase is its standard cosine parametrization. Governing-law hypotheses below state how these scalar objects represent this rotational system
\end{definition}
\begin{proof}
  This is the declared model definition of Torsion Pendulum Setup.
\end{proof}
\begin{definition}[Matches Problem Data]
  \label{def:physics:phyx-mini-0278:phyxminiproblems-problemphyxmini0278-matchesproblemdata}
  \lean{PhyXMiniProblems.ProblemPhyXMini0278.MatchesProblemData}
  \uses{def:physics:phyx-mini-0278:phyxminiproblems-problemphyxmini0278-timeinseconds, def:physics:phyx-mini-0278:phyxminiproblems-problemphyxmini0278-torquemagnitudeinnewtonmeters, def:physics:phyx-mini-0278:phyxminiproblems-problemphyxmini0278-angularvelocityinradianspersecond, def:physics:phyx-mini-0278:phyxminiproblems-problemphyxmini0278-torsionpendulumsetup}
  Numerical and qualitative facts stated in the problem text. In particular, the disk is released from rest at 0.200 rad, tau\_s = 0.004 N m, and t\_s = 0.40 s. No maximum-speed value or answer choice occurs here
\end{definition}
\begin{proof}
  This is the declared model definition of Matches Problem Data.
\end{proof}
\begin{definition}[Matches Supplied Figure]
  \label{def:physics:phyx-mini-0278:phyxminiproblems-problemphyxmini0278-matchessuppliedfigure}
  \lean{PhyXMiniProblems.ProblemPhyXMini0278.MatchesSuppliedFigure}
  \uses{def:physics:phyx-mini-0278:phyxminiproblems-problemphyxmini0278-timeinseconds, def:physics:phyx-mini-0278:phyxminiproblems-problemphyxmini0278-torquemagnitudeinnewtonmeters, def:physics:phyx-mini-0278:phyxminiproblems-problemphyxmini0278-torsionpendulumsetup}
  Quantitative and qualitative evidence transcribed from the primary bitmap. The lower graph begins at a positive maximum, reaches the negative maximum at half of t\_s, and returns to the next positive maximum at t\_s; thus the marked horizontal scale is one oscillation period. These are data readouts, not assumptions about the requested maximum angular speed
\end{definition}
\begin{proof}
  This is the declared model definition of Matches Supplied Figure.
\end{proof}
\begin{definition}[Has Physical Torsion Pendulum Parameters]
  \label{def:physics:phyx-mini-0278:phyxminiproblems-problemphyxmini0278-hasphysicaltorsionpendulumparameters}
  \lean{PhyXMiniProblems.ProblemPhyXMini0278.HasPhysicalTorsionPendulumParameters}
  \uses{def:physics:phyx-mini-0278:phyxminiproblems-problemphyxmini0278-timeinseconds, def:physics:phyx-mini-0278:phyxminiproblems-problemphyxmini0278-momentofinertiainkilogrammetersquared, def:physics:phyx-mini-0278:phyxminiproblems-problemphyxmini0278-torquemagnitudeinnewtonmeters, def:physics:phyx-mini-0278:phyxminiproblems-problemphyxmini0278-torsionconstantinnewtonmetersperradian, def:physics:phyx-mini-0278:phyxminiproblems-problemphyxmini0278-torsionpendulumsetup}
  Positivity and nondegeneracy of the physical readouts
\end{definition}
\begin{proof}
  This is the declared model definition of Has Physical Torsion Pendulum Parameters.
\end{proof}
\begin{definition}[Satisfies Ideal Torsion Pendulum Laws]
  \label{def:physics:phyx-mini-0278:phyxminiproblems-problemphyxmini0278-satisfiesidealtorsionpendulumlaws}
  \lean{PhyXMiniProblems.ProblemPhyXMini0278.SatisfiesIdealTorsionPendulumLaws}
  \uses{def:physics:phyx-mini-0278:phyxminiproblems-problemphyxmini0278-momentofinertiainkilogrammetersquared, def:physics:phyx-mini-0278:phyxminiproblems-problemphyxmini0278-torquemagnitudeinnewtonmeters, def:physics:phyx-mini-0278:phyxminiproblems-problemphyxmini0278-torsionconstantinnewtonmetersperradian, def:physics:phyx-mini-0278:phyxminiproblems-problemphyxmini0278-angularvelocityinradianspersecond, def:physics:phyx-mini-0278:phyxminiproblems-problemphyxmini0278-torsionpendulumsetup}
  Governing laws of the ideal undamped torsion pendulum. The first two fields identify Physlib's positive scalar oscillator parameters with rotational inertia and torsion constant. The third field is the linear torque-magnitude law shown by the upper graph. The remaining fields express the zero-phase cosine solution and identify angular velocity with dtheta/dt. They contain no assertion about a maximum and no displayed answer value
\end{definition}
\begin{proof}
  This is the declared model definition of Satisfies Ideal Torsion Pendulum Laws.
\end{proof}
\begin{definition}[Answer Choice]
  \label{def:physics:phyx-mini-0278:phyxminiproblems-problemphyxmini0278-answerchoice}
  \lean{PhyXMiniProblems.ProblemPhyXMini0278.AnswerChoice}
  Labels of the four numerical choices printed with the problem
\end{definition}
\begin{proof}
  This is the declared model definition of Answer Choice.
\end{proof}
\begin{definition}[displayed Angular Speed Radians Per Second]
  \label{def:physics:phyx-mini-0278:phyxminiproblems-problemphyxmini0278-displayedangularspeedradianspersecond}
  \lean{PhyXMiniProblems.ProblemPhyXMini0278.displayedAngularSpeedRadiansPerSecond}
  \uses{def:physics:phyx-mini-0278:phyxminiproblems-problemphyxmini0278-answerchoice}
  Displayed angular-speed values in radians per second
\end{definition}
\begin{proof}
  This is the declared model definition of displayed Angular Speed Radians Per Second.
\end{proof}
\begin{definition}[recorded Dataset Answer]
  \label{def:physics:phyx-mini-0278:phyxminiproblems-problemphyxmini0278-recordeddatasetanswer}
  \lean{PhyXMiniProblems.ProblemPhyXMini0278.recordedDatasetAnswer}
  \uses{def:physics:phyx-mini-0278:phyxminiproblems-problemphyxmini0278-answerchoice}
  The answer label recorded in the dataset metadata
\end{definition}
\begin{proof}
  This is the declared model definition of recorded Dataset Answer.
\end{proof}
\begin{definition}[Matches Answer Choice]
  \label{def:physics:phyx-mini-0278:phyxminiproblems-problemphyxmini0278-matchesanswerchoice}
  \lean{PhyXMiniProblems.ProblemPhyXMini0278.MatchesAnswerChoice}
  \uses{def:physics:phyx-mini-0278:phyxminiproblems-problemphyxmini0278-answerchoice, def:physics:phyx-mini-0278:phyxminiproblems-problemphyxmini0278-displayedangularspeedradianspersecond}
  Agreement with a value displayed to the nearest 0.01 rad/s
\end{definition}
\begin{proof}
  This is the declared model definition of Matches Answer Choice.
\end{proof}
% --- Archon declaration topology end ---
% --- Archon physics formalization source end ---
